\section{Collaboration}
\label{sec:collaboration}

Collaboration is important throughout this project because the different
activities are connected. Data discovery affects extraction, while the
extracted data influences the ERD, normalization, database implementation,
and competency queries. Although each member has a main responsibility,
the group works together whenever one task affects another part of the
project.

\subsection{Intra-Level Collaboration}
\label{subsec:intra-level-collaboration}

Intra-level collaboration refers to the cooperation among the four members
of Group 07. The members divide the work according to their main
responsibilities. However, these responsibilities are not treated as
completely separate tasks. Members regularly help one another, review
shared outputs, and discuss any problem that affects the project.

\subsubsection{Member Roles and Shared Responsibilities}

The main responsibilities of the group members are shown in
Table~\ref{tab:collaboration-roles}.

\begin{table}[htbp]
    \centering
    \caption{Main responsibilities of the group members}
    \label{tab:collaboration-roles}
    \small
    \renewcommand{\arraystretch}{1.2}

    \begin{tabularx}{\textwidth}{
        @{}
        >{\raggedright\arraybackslash}p{0.26\textwidth}
        >{\raggedright\arraybackslash}X
        @{}
    }
        \toprule
        \textbf{Member} & \textbf{Main responsibilities} \\
        \midrule

        Sheikh Tawsif Bin Aftab &
        Works as the Team Lead and Scrum Master. He manages the overall progress
        of the group, maintains the Trello board, distributes tasks, records regular
        work updates, participates in data extraction, implements the database, and
        prepares competency queries. \\

Syed Mehraz Monwar &
Participates in data discovery and contributes to documentation, report
structure, drawing, review, and verification. \\

Sanjida Mahi &
Reviews source material, prepares and revises the ERD, verifies the model after
each major change, and works on normalization from the initial relations
through BCNF. \\

Safwan Ahmed Rayed &
Participates in data discovery and works with Sheikh Tawsif Bin
Aftab during data extraction. He also helps with ERD preparation,
verification, normalization, and documentation. \\

        \bottomrule
    \end{tabularx}
\end{table}

The responsibilities in Table~\ref{tab:collaboration-roles} identify the
main areas of work rather than setting strict boundaries. For example,
Sanjida Mahi mainly works on the ERD and normalization, but Safwan Ahmed
Rayed also supports these activities. Similarly, Sheikh Tawsif Bin Aftab
and Safwan Ahmed Rayed work together during extraction, while Syed Mehraz
Monwar contributes to data discovery with the other members. In this way,
the team completes connected tasks through shared effort.



\subsubsection{Task Coordination and Progress Updates}

The group uses Trello to record tasks and monitor their progress. As the
Scrum Master, Sheikh Tawsif Bin Aftab manages the Trello board. He adds and
distributes tasks, updates their status, and follows the overall progress of
the team. This allows the members to identify completed, ongoing, and
remaining activities.

The team also uses shared working materials for source files, extracted
data, ERD working files, database files, and report sections. When a member
updates an important output, the other members review it before it is
accepted as part of the project.

\subsubsection{Internal Review and Decision Making}

Review takes place throughout the project instead of only at the end.
The members check the collected data, extracted files, ERD revisions,
relations, primary keys, and normalization results. This is necessary
because an error in one stage can create further problems during database
implementation.

The ERD receives repeated review as the team gains a better understanding of
the source data. Sanjida Mahi prepares and verifies the approved final drawing,
with support from Safwan Ahmed Rayed and the other team members. The relations
are also checked step by step until the final database satisfies BCNF.

\subsubsection{Team Meetings and Problem Solving}

The group holds meetings whenever a problem arises or when an activity
requires cooperation from the whole team. Therefore, the meetings are
arranged according to the needs of the project rather than following only
a fixed schedule.

During these meetings, the members report their progress and discuss
problems found in the source files, extraction process, ERD, normalization,
and database implementation. Possible solutions are compared before the
team reaches a shared decision. The responsible member then updates the
related work, and the result is reviewed again when necessary.

This process allows the members to complete their individual
responsibilities while remaining involved in the overall development of
the project.

\subsection{Inter-Level Collaboration}
\label{subsec:inter-level-collaboration}

Inter-level collaboration refers to coordination between Group 07 and the
wider MODELBD Scrum-of-Scrums structure. Common milestones, review
expectations, and regular academic guidance affect the group's work, but the
available evidence does not establish a direct exchange of datasets, database
relations, ERD components, or implementation responsibilities with another
group. This distinction keeps the collaboration description consistent with
the work that can be verified.

\subsection{Competency Questions as a Collaboration Tool}
\label{subsec:competency-collaboration}

Competency questions give the team a common way to check whether modelling,
normalization, implementation and deployment agree. A modeller can state that water
quality connects to rivers, but the claim is accepted only when the implementer
runs the join and the reviewer checks its count. The same questions are saved
as SQL and automated tests, so they can be repeated after a source or schema
change.

Figure~\ref{fig:competency-collaboration} shows how the competency questions
connect modelling, implementation and final review.

\begin{figure}[H]
\centering
\begin{tikzpicture}[
  role/.style={draw=cublue,rounded corners,fill=cugrey,text width=2.6cm,
    minimum height=0.95cm,align=center,font=\small},
  evidence/.style={draw=cuteal,rounded corners,fill=white,text width=3.0cm,
    minimum height=0.95cm,align=center,font=\small},
  arr/.style={-{Stealth[length=2mm]},thick,draw=cuteal},
  node distance=0.55cm
]
\node[role] (model) {Model and normalize};
\node[role,right=of model] (implement) {Implement and load};
\node[role,right=of implement] (review) {Review and deploy};
\node[evidence,below=0.75cm of implement] (cq) {Shared competency questions,
expected results and timing};
\draw[arr] (model)--(implement);
\draw[arr] (implement)--(review);
\draw[arr] (model)--(cq);
\draw[arr] (implement)--(cq);
\draw[arr] (review)--(cq);
\end{tikzpicture}
\caption{Competency questions connect intra-level work products.}
\label{fig:competency-collaboration}
\end{figure}
